\scanpage{27}
\mspage{\handwritten{19}}
% The source is upside down; read after a 180-degree rotation.
% The entire draft is cancelled by long diagonal strokes.
\margin{\handwritten{\textcircled{$\unsure{\psi|\Psi}$}}}

\deleted{2.1.1. On peut dire encore qu'une structure de biextension sur le $G_{P\times Q}$-torseur $E$ est définie par deux lois de composition non partout définies $\handwritten{\phi}(x,y)$ et $\handwritten{\unsure{\mu|\psi}}(x,y)$ sur $E$, la première étant définie lorsque \deleted{\unsure{$j($}} $\mathrm{pr}_2j(x)=\mathrm{pr}_2j(y)$, la deuxième lorsque $\mathrm{pr}_1j(x)=\mathrm{pr}_1j(y)$ \deleted{(où} (où $j:E\handwritten{\longrightarrow}P\times Q$ est le morphisme structural du torseur). Les axiomes, en plus des associativités du type}
\[
\deleted{\deleted{\handwritten{\phi}(gx,y)\unsure{=}g\handwritten{\phi}(x,y)\unsure{,}\quad\handwritten{\phi}(x,gy)\unsure{=}g\handwritten{\phi}(x,y)\quad\handwritten{\phi}(}}
\]
\[
\deleted{\handwritten{\phi}(gx,y)=\handwritten{\phi}(x,gy)=g\handwritten{\phi}(x,y)\qquad(x,y\handwritten{\in}E(S),\ g\handwritten{\in}G(S))}
\tag*{(2.1.1.1)}
\]

\deleted{et des relations analogues en $\handwritten{\mu}$, s'expriment alors par l\deleted{a}'associativité et la commutativité de \unsure{ces|ses} lois de composition partielles, et par la relation}
\[
\deleted{\handwritten{\phi}(\handwritten{\mu}(x,y),\handwritten{\mu}(z,t))=\handwritten{\mu}(\handwritten{\phi}(x,z),\handwritten{\phi}(y,t))}
\tag*{(2.1.1.2)}
\]
\deleted{pour $x,y,z,t$ respectivement dans $E_{p,q}(S),E_{p,q'}(S),E_{p',q}(S),E_{p',q'}(S)$.}

\deleted{2.2. \inserted{\handwritten{\uline{Sections unité}.}} Il y a lieu de considérer la section unité de $\mathbb{E}\deleted{\unsure{t|E}}$ considéré comme extension de $Q_P$ par $G_P$,}
\[
\deleted{e_{E/P}\handwritten{\in\Gamma}(E/P)=\mathrm{Hom}_P(P,E)\ ,}
\tag*{(2.2.1)}
\]
\deleted{et de façon symétrique la section unité de $E$ considéré comme extension de $P_Q$ par $G_Q$,}
\[
\deleted{e_{E/Q}\handwritten{\in\Gamma}(E/Q)=\mathrm{Hom}_Q(Q,E)\ .}
\tag*{(2.2.2)}
\]

\deleted{En particulier, \uline{le torseur} $E$ \uline{au dessus de} $P\times Q$ (de \unsure{G|g}roupe $G_{P\times Q}$) \uline{est canoniquement trivialisé au dessus de} $1_P\times Q$ \uline{par} $e_{E/P}$, \uline{et au dessus de} $P\times1_Q$ \uline{par} $e_{E/Q}$. \inserted{De plus,} \uline{Ces deux trivialisations coïncident au dessus de la section unité} $1_P\times1_Q$ de $P\times Q$, et définissent\inserted{\handwritten{donc}} une même section}
\[
\deleted{e_{E\deleted{\unsure{/S|/P|/Q}}}\handwritten{\in\Gamma}(E\handwritten{)}\deleted{\unsure{/S)|/P)|/Q)}}\ ,}
\tag*{(2.2.3)}
\]
\deleted{qu'on appellera par abus de langage la \uline{section unité} (ou \uline{section neutre}) de $E$. Pour prouver l'égalité annoncée\deleted{\unsure{et une conséquence}} des deux sections $e,e'$ de $E$ déduites respectivement de $e_{E/P}$ et $e_{E/Q}$, on note}
% archive sheet number: 26
